\begin{figure}[tbp]
\centering
\begin{tikzpicture}[
  phase/.style={draw=black!65,rounded corners=2pt,align=center,
    minimum height=1.25cm,text width=2.8cm,font=\small,inner sep=4pt},
  >={Stealth[length=2mm]}
]
\node[phase,fill=blue!5] (reset) {Reset\\\(t\in[2,5]\)\\\(v'_\ell=0\)};
\node[phase,fill=blue!5,right=0.35cm of reset] (prep)
  {Preparation\\\(t\in(8,12)\)\\add \(U_\ell(t,z)\)};
\node[phase,fill=orange!8,right=0.35cm of prep] (slow)
  {Possible delay\\\(t\in(48,52)\)\\speed set by \(z\)};
\node[phase,fill=blue!5,right=0.35cm of slow] (read)
  {Readout\\\(t\in(74,80)\)\\add aligned terms to \(z^2\)};
\draw[->] (reset)--(prep);
\draw[->] (prep)--(slow);
\draw[->] (slow)--(read);
\draw[|-|,black!85] ([yshift=-0.45cm]prep.south west)
  -- node[below,font=\small,align=center]
  {\(a(t)=L\) throughout preparation, delay and readout}
  ([yshift=-0.45cm]read.south east);
\end{tikzpicture}
\caption{Clock intervals during a complete passage, shown schematically
and not to scale. Values refer to the clock before each update.
All registers reset before preparation. Each channel is read at most
once; for a channel that is read, its preparation signals give a bounded
contribution at that readout.}
\label{fig:clock}
\end{figure}
